\documentclass[11pt,letterpaper]{article}
\usepackage[letterpaper,margin=0.22in]{geometry}
\usepackage[utf8]{inputenc}
\usepackage{mdwlist}
\usepackage{charter}
\usepackage[T1]{fontenc}
\usepackage{textcomp}
\usepackage{enumitem}
\pagestyle{empty}
\setlength{\tabcolsep}{0em}

\newenvironment{indentsection}[1]%
{\begin{list}{}%
{\setlength{\leftmargin}{#1}}%
     \item[]
}
{\end{list}}
\newenvironment{unindentsection}[1]%
{\begin{list}{}%
{\setlength{\leftmargin}{-0.5#1}}%
\item[]%
}
{\end{list}}
\newcommand{\headerrow}[2]
{\begin{tabular*}{\linewidth}{l@{\extracolsep{\fill}}r}
#1 &
#2 \\
\end{tabular*}}
\newcommand{\CPP}
{C\nolinebreak[4]\hspace{-.05em}\raisebox{.22ex}{\footnotesize\bf ++}}
\begin{document}

\begin{center}
	{\LARGE \textbf{Raj Reddy}}\\
	\vspace{0.05cm}
	(352) 530-3397
    \hfill rajreddy23@outlook.com
    \hfill github.com/rjrddy
    \hfill linkedin.com/in/raj-reddy-1
    \hfill raj-reddy.com
\end{center}
\hrule
\vspace{\stretch{1}}

\vspace{-1em}
\subsection*{\Large Education}
\renewcommand\labelitemi{}
\begin{itemize}[leftmargin=1em]
	\parskip=0.1em
	\item
	      \headerrow
	      {\textbf{University of Utah}}
	      {\textbf{Salt Lake City, UT}}
	      \headerrow
	      {\emph{B.S. Computer Science}}
	      {\emph{August 2020 -- May 2025}}
	      \item \textbf{Relevant Coursework:} Computer Systems, Machine Learning, Computer Graphics, Algorithms, Software Practice I \& II, Database Systems, Computer Networks, Foundations of Data Analysis, Models of Computation, and Linear Algebra.
\end{itemize}
\hrule
\vspace{\stretch{1}}

\vspace{-1em}
\subsection*{\Large Skills}
\hyphenpenalty=1000
\begin{itemize}[leftmargin=1em]
    \parskip=0.1em
	\item \textbf{Languages:} Python, TypeScript, C++, SQL, Java
	\item \textbf{Software:} Pandas, NumPy, React, Node.js, AWS, Azure, Docker, Next.js, .NET, REST, Django, MongoDB, Linux
\end{itemize}
\hrule
\vspace{\stretch{1}}

\vspace{-1em}
\subsection*{\Large Experience}
\renewcommand\labelitemi{}
\renewcommand\labelitemii{$\bullet$}
\begin{itemize}[leftmargin=1em]
	\parskip=0.1em

	\item
	      \headerrow
	      {\textbf{L3Harris Technologies}}
	      {\textbf{Dallas, TX}}
	      \headerrow
	      {\emph{Associate Software Engineer}}
	      {\emph{June 2025 -- Present}}
	      \begin{itemize*}
	      	\item Developed real-time signal processing algorithms in Python and \CPP{} on embedded systems, applying DSP techniques for defense sensor data streams.
	      	\item Built microservices for high-frequency data ingestion and storage, leveraging Podman and serverless infrastructure on AWS.
	      	\item Automated CI/CD pipelines with static code analysis and integration testing.
	      \end{itemize*}

	\item
	      \headerrow
	      {\textbf{Doxy.me}}
	      {\textbf{Charleston, SC}}
	      \headerrow
	      {\emph{Software Engineer}}
	      {\emph{August 2024 -- May 2025}}
	      \begin{itemize*}
	      	\item Improved diagnostic throughput by building end-to-end ML inference pipelines in Python and TypeScript, orchestrating model execution with PyTorch and Llama over unstructured clinical data ingested from WebRTC video sessions, with results persisted to downstream analytics on AWS SageMaker and Fargate.
	      	\item Reduced data processing latency by designing and optimizing ETL workflows that aggregated and normalized high-volume clinical records across multiple input schemas, enabling reliable daily pipeline execution at scale.
	      	\item Increased release cadence by owning CI/CD configuration, cloud infrastructure automation, and observability tooling on AWS, giving the team confidence to ship multiple features per week without operational overhead.
	      \end{itemize*}

	\item
	      \headerrow
	      {\textbf{University of Utah}}
	      {\textbf{Salt Lake City, UT}}
	      \headerrow
	      {\emph{Undergraduate Research Assistant -- FuTURES Lab}}
	      {\emph{May 2024 -- Jun 2025}}
	      \begin{itemize*}
	      	\item Analyzed and optimized large-scale datasets, enhancing software configuration testing with gcov and CMake, increasing code coverage by 30\% on real-world APIs (Libpng, PyTorch).
	      	\item Expanded OSS-Fuzz testing coverage using compile-time options to improve the robustness of full-stack libraries.
	      \end{itemize*}
\end{itemize}
\hrule
\vspace{\stretch{1}}

\vspace{-1em}
\subsection*{\Large Projects}
\hyphenpenalty=1000
\begin{itemize}[leftmargin=1em,noitemsep]
\item \textbf{Telehealth Mental Health Diagnosis Tool (Doxy.me):} Built an AI system that analyzes visual, speech, and tonal cues from telehealth sessions using PyTorch and Llama, delivering structured mental health insights to treating physicians via SageMaker and Fargate inference pipelines. Improved diagnostic accuracy using reinforcement learning and WebRTC video integration.
\item \textbf{University Course Planning/Review Platform} \emph{(Class Best):} Built a multi-service data platform enabling student schedule matching, course overlap detection, and behavioral insights. Developed ETL pipelines in Python to process thousands of course records, user preferences, and time-series activity logs. Built a responsive frontend in React + Tailwind with interactive dashboards, heatmaps, and calendar visualizations.
\item \textbf{Configuration Fuzzer:} Developed a configuration fuzzing tool for OSS-Fuzz, automating build generation to identify critical compile-time configurations and improve code coverage, uncovering previously untested code paths and increasing bug detection effectiveness.
\end{itemize}
\vspace{\stretch{1}}
\end{document}
